% Generated draft: requires mathematical and source review.
\documentclass[11pt]{article}
\usepackage[margin=25mm]{geometry}
\usepackage{fontspec}
\setmainfont{Latin Modern Roman}
\usepackage{amsmath,amssymb,mathtools,unicode-math}
\setmathfont{Latin Modern Math}
\usepackage{enumitem,needspace,xurl,hyperref}
\setlength{\parindent}{0pt}
\setlength{\parskip}{5pt}
\setlength{\emergencystretch}{3em}
\newcommand{\N}{\mathbb{N}}
\newcommand{\Z}{\mathbb{Z}}
\newcommand{\Q}{\mathbb{Q}}
\newcommand{\R}{\mathbb{R}}
\newcommand{\C}{\mathbb{C}}
\newcommand{\F}{\mathbb{F}}
\newcommand{\A}{\mathbb{A}}
\newcommand{\PP}{\mathbb{P}}
\newcommand{\E}{\mathbb{E}}
\newcommand{\eps}{\varepsilon}
\newcommand{\dd}{\,\mathrm{d}}
\begin{document}
\section*{UnsolvedMath ID: 30004756}
\textbf{Source problem code:} catalogue entry
\subsection{Definitions and mathematical statement}
The supplied record concerns dimensions \(n\in\N\), \(n\geq5\), and the formal equation
\[
\det D^2u=1+\sum_{i\in I}a_i\delta_{p_i},
\]
where \(u\) is convex, \(p_i\) are mass locations, \(a_i\) are mass sizes, and \(\delta_p\) is the unit Dirac measure at \(p\). The index set \(I\), domain, admissible masses, boundary conditions, and normalization are not specified in the supplied record.

A standard measure interpretation can be described conditionally as follows. For a specified open convex domain \(\Omega\subseteq\R^n\) and a finite-valued convex function \(u:\Omega\to\R\), define
\[
\partial u(x)=\{q\in\R^n:u(y)\geq u(x)+q\cdot(y-x)\text{ for every }y\in\Omega\},
\qquad
\partial u(E)=\bigcup_{x\in E}\partial u(x).
\]
The Aleksandrov Monge--Amp\`ere measure is
\[
\operatorname{MA}_u(E)=\mathcal L^n(\partial u(E))
\]
for Borel sets \(E\subseteq\Omega\), where \(\mathcal L^n\) is Lebesgue measure. With this interpretation, the displayed equation means
\[
\operatorname{MA}_u=\mathcal L^n\!\restriction_\Omega+\sum_{i\in I}a_i\delta_{p_i},
\]
with \(p_i\in\Omega\) and with the right-hand side required to define a nonnegative locally finite measure. Here the scalar \(1\) in the formal equation represents Lebesgue density one. This interpretation is supplied to explain the notation; the record does not explicitly identify its solution notion or assert that this domain framework is the intended one. It does not specify whether \(I\) is finite, whether the points are distinct, or what sign restrictions are imposed on individual masses.

\textbf{Requested mathematical targets.} For the intended class of convex solutions and its unspecified parameter space, the source asks:
\begin{enumerate}
\item For dimensions \(n\geq5\), determine whether singular solutions analogous to the lower-dimensional examples referenced in the report exist. The examples and the precise singular structures to be reproduced are not included in the supplied record.
\item Determine whether the relevant singularities persist under small perturbations of the locations \(p_i\) and masses \(a_i\). The admissible perturbations, their topology, the choice or existence of corresponding solutions, and the meaning of persistence are unspecified.
\item Determine the dimension, regularity, connectivity, and other geometric properties of the locus of singular solutions in the source's ``solution orbifold,'' including algebro-geometric conditions that exclude singular structures. The solution orbifold, its equivalence relation, and its geometric structure are not defined in the supplied record.
\end{enumerate}
A spatial singular set of one solution and a locus of singular solutions in a parameter or moduli space are different objects. The original statement explicitly requests the latter, whereas the supplied literature assessment refers to regularity and connectivity of ``the singular set.'' The record does not resolve whether additional spatial singular-set questions are intended, nor does it specify a regularity criterion defining singularity. Consequently these targets cannot be converted into a uniquely quantified conjecture without additional source information.

The supplied status is \emph{partially solved}. Its dated curation reports that higher-dimensional polyhedral singular structures and stability under perturbation have been proved, while the extracted regularity and connectivity questions for the singular set remain unsettled in general. These are supplied literature claims, not independently verified results or a new assessment of which formulations remain open.
\subsection{Short English statement}
For convex Monge--Amp\`ere solutions with Lebesgue density one and Dirac masses, the source asks about singular examples in dimensions \(n\geq5\), persistence under perturbing the masses, and the geometry of the singular-solution locus in a purported solution orbifold. The domain, solution class, singularity criterion, perturbation conditions, and orbifold are unspecified. The supplied dated curation labels the problem partially solved and reports progress on polyhedral structures and stability; its precise scope requires review.
\subsection{Sources}
\begin{itemize}
\item Local UnsolvedMath registry, supplied complete source record: permanent ID \(30004756\), title \emph{Geometry of Singular Monge--Amp\`ere Solution Loci}, collection \emph{Oberwolfach Reports Open Problems}, exact source code \texttt{catalogue entry}. No individual proposer is supplied; the proposed year is recorded as 2021.
\item Supplied report citation: \emph{Partial Differential Equations (2022)}, Oberwolfach Reports, pp.\ 1859--1914, DOI \texttt{10.4171/owr/2021/35}. \url{https://doi.org/10.4171/owr/2021/35}. The record describes Oberwolfach Reports as documenting mathematical workshops and their open-problem discussions. The citation's 2022 label and the DOI's 2021 component are retained as supplied.
\item Supplied literature source, labelled ``Mooney--Rakshit'': \url{https://arxiv.org/abs/2204.11365}. No title or theorem-level attribution is supplied.
\item Supplied literature source, labelled ``Mooney'': \url{https://par.nsf.gov/servlets/purl/10213707}. No title or theorem-level attribution is supplied.
\item Literature assessment date: \texttt{2026-08-22}. This is supplied dated curation, not independently verified. It labels the problem partially solved and reports proofs of higher-dimensional polyhedral singular structures and stability under perturbation, while stating that the extracted regularity and connectivity questions for the singular set are not settled in general. The background calls this a dated literature triage derived from the supplied OWR curation and requests independent verification of time-sensitive status and cited claims.
\item The registry labels the statement \texttt{corrected\_verified} and states that the official report was checked directly and that the clean statement preserves the complete target while restoring definitions and removing extraction artifacts or source-dependent wording. This is a supplied verification claim; no independent source consultation is asserted here.
\end{itemize}
\end{document}
